% ---------------------------------------------------------------------------
% Conclusion
%
% Answer the title in the first sentence, with a number: how far HL goes, and
% what stops it. Then the budget condition under which that answer holds, and
% what would change it.
%
% THE ANSWER MUST CARRY THE SCOPE, as of 2026-08-19. The title now reads "in Adversarial
% Games". So the first sentence says how far HL goes ON GAMES -- not how far policy-level
% continual learning goes. One qualifying clause, no more: an unscoped takeaway here would
% undo the limitation, and a paragraph about scope would repeat Section 1.
%
% Keep it short. Do not restate the contribution list from the introduction.
% ---------------------------------------------------------------------------

\section{Conclusion}

We define Adversarial Heuristic Learning (AHL), in which a coding agent reads
decision-level game feedback and revises an executable symbolic policy. We introduce
AHL-Arena, a reproducible benchmark built from a diverse archive of real adversarial
games, and evaluate AHL against human programs and reinforcement-learning baselines
under matched interaction budgets. Our results show when the readability and rich
feedback channel of program policies outweigh their maintenance cost, and when a
learned policy remains the better choice; together, they establish adversarial games
as a concrete setting for measuring how far self-evolving coding agents can carry
heuristic programs.
